Joint-Embedding Predictive Architectures (JEPAs) have become a leading
recipe for latent world models: an encoder maps each observation to a
representation and a predictor forecasts future representations. Yet
these models are \emph{memoryless} in time --- the encoder
sees one frame at a time and the predictor attends only over a short
fixed window --- so they cannot represent how information at
different temporal distances shapes the dynamics of the latent space. We
introduce a minimal, mathematically transparent remedy: a
\textbf{two-timescale exponential memory} that augments the predictor
with a \emph{fast} (short-term) and a \emph{slow} (long-term)
exponential-moving-average (EMA) bank over the latent stream, injected
through zero-initialized projections so training begins exactly at the
memoryless baseline. The mechanism is the simplest diagonal linear
state-space model; it adds two interpretable scalars whose closed-form
effective horizon is \(\tau=-1/\ln(1-\alpha)\), and it keeps the host model's
two-term loss intact. On a controlled suite of partially-observable,
memory-stressing environments built on the recent LeWorldModel, we show
that (i) the memory horizon must \emph{match} the task's cue-to-decision
gap --- a short bank bridges short gaps and only the long bank
bridges long ones; (ii) the matched timescale carries information to the
\emph{decision}, lifting cue-decodability from the model's own
prediction from chance to \(0.84\) on long-horizon tasks; and (iii)
memory extends the usable horizon far beyond the predictor's window,
degrading gracefully where the memoryless baseline cliffs. A
fully-observable control confirms the gains are memory-specific. Across
5 seeds we further show the matched timescale \textbf{causally} drives
both the prediction (a counterfactual memory-swap flips the prediction
to an injected cue) and \textbf{downstream closed-loop control}
(test-time memory ablation collapses success to chance); a fixed
\textbf{log-spaced multi-bank} captures the full range of horizons
\emph{without tuning} and outperforms a learned GRU and a long-context
predictor. We argue that an explicit, controllable multi-timescale
memory --- rather than the hierarchy/sub-goal direction the
field currently favors --- is a foundational primitive for
studying memory in JEPA latent dynamics, and we contribute a reusable
\emph{availability-vs-usage} measurement protocol.

\begin{center}\rule{0.5\linewidth}{0.5pt}\end{center}
